\documentclass[11pt]{article}
\usepackage[margin=0.85in]{geometry}
\usepackage{amsmath,amssymb,amsthm,booktabs,array,longtable}
\usepackage[T1]{fontenc}
\usepackage[colorlinks=true,urlcolor=blue,citecolor=blue,linkcolor=blue]{hyperref}
\newtheorem{theorem}{Theorem}
\newtheorem{lemma}[theorem]{Lemma}
\newtheorem{corollary}[theorem]{Corollary}
\newcommand{\A}{A_{u,v}}
\title{Proof of Bala's totient-period conjectures:\\
OEIS A258899 and A258903, COMMENTS\\
(3 and 8 August 2025)}
\author{Haobo Ma\\\small ChronoAI Pte Ltd, Singapore
\and Wenlin Zhang\\\small National University of Singapore\\\small ChronoAI Pte Ltd, Singapore}
\date{3 November 2026}
\begin{document}
\maketitle
\begin{abstract}
We prove the statements labelled ``Conjecture'' in Peter Bala's COMMENTS
of 3 August 2025 in OEIS A258899 and 8 August 2025 in A258903
\cite{a258899,a258903}: modulo every positive integer $m$, each sequence
is eventually periodic with least eventual period dividing $\varphi(m)$.
For both sequences, $\varphi(m)$ is a period from $n=2m+4$ onward.
The same factorial-divisibility argument proves Bala's constant-tail
conjecture in A121860, COMMENTS, 8 August 2025 \cite{a121860}.
More generally, the divisor sum with weights $u^d v^{n/d}$ reduces modulo
$m$ to $uv^n+vu^n$ for $n\geq2m+4$; the Carmichael exponent is also a period.
For prime moduli $p$, the onset improves to $n=p$, proving the numerical
onsets in the two entries. Further corollaries cover the directly linked
factorial sums and correct a three-consecutive-residues observation
in the history of A121860 by a counterexample modulo $89$.
\end{abstract}

\section{The sequences and the result}
Paul D. Hanna's entries A258899 and A258903 have offset zero and initial
value $1$ \cite{a258899,a258903}. Their FORMULA fields, contributed by
Bala on 4 and 8 August 2025 respectively, give, for $n\geq1$,
\begin{align*}
 a(n)&=\sum_{d\mid n}2^d\frac{n!}{d!(n/d)!},&
 b(n)&=\sum_{d\mid n}2^d3^{n/d}\frac{n!}{d!(n/d)!}.
\end{align*}
The linked entry A121860, originally contributed by Vladeta Jovovic,
has offset one and is
$c(n)=\sum_{d\mid n} n!/(d!(n/d)!)$ for $n\geq1$ \cite{a121860}.

For integers $u,v$, set $\A(0)=1$ and
\begin{equation}\label{eq:def}
 \A(n)=\sum_{d\mid n}u^dv^{n/d}Q(n,d),\qquad
 Q(n,d)=\frac{n!}{d!(n/d)!}\quad(n\geq1).
\end{equation}
These quotients are integers. If $d=1$ or $d=n$ they equal $1$.
Otherwise $d,e=n/d\geq2$ and $d+e\leq de=n$, so
$Q(n,d)=\binom{n}{d,e,n-d-e}(n-d-e)!$.
Thus $a=A_{2,1}$, $b=A_{2,3}$, and $c(n)=A_{1,1}(n)$ for $n\geq1$.
Write $\varphi$ for Euler's totient, with $\varphi(1)=1$.

\begin{theorem}\label{thm:main}
For $u,v\in\mathbb Z$, $m\geq1$, and $n\geq2m+4$,
\begin{equation}\label{eq:tail}
 \A(n)\equiv uv^n+vu^n\pmod m,
\end{equation}
and consequently
\begin{equation}\label{eq:period}
 \A(n+\varphi(m))\equiv\A(n)\pmod m.
\end{equation}
\end{theorem}

\section{Factorial divisibility and powers}
\begin{lemma}\label{lem:factorial}
If $m\geq1$, $d,e\geq2$, $n=de$, and $n\geq2m+4$, then
$m!\mid n!/(d!e!)$.
\end{lemma}
\begin{proof}
The inequality $(d-2)(e-2)\geq0$ gives
$2(d+e)\leq n+4$, and hence $d+e+m\leq n$.
Put $r=n-d-e-m\geq0$. The multinomial coefficient
$n!/(d!e!m!r!)$ is an integer, so
\[
 \frac{n!}{d!e!}=m!r!\binom{n}{d,e,m,r}
\]
is divisible by $m!$.
\end{proof}

\begin{lemma}\label{lem:powers}
For every integer $z$, $m\geq1$, and $n\geq m$,
$z^{n+\varphi(m)}\equiv z^n\pmod m$.
\end{lemma}
\begin{proof}
The case $m=1$ is immediate. For each prime power $p^e\Vert m$,
if $p\mid z$, both powers vanish modulo $p^e$ because $n\geq m\geq e$.
If $p\nmid z$, multiplication by $z$ permutes the units modulo $p^e$.
Taking the product of all units and cancelling that product gives
$z^{\varphi(p^e)}\equiv1\pmod{p^e}$.
The formula $\varphi(m)=\prod_{q^f\Vert m}q^{f-1}(q-1)$ shows
$\varphi(p^e)\mid\varphi(m)$, so the two powers are again congruent.
The Chinese remainder theorem combines these congruences modulo $m$.
\end{proof}

\begin{proof}[Proof of Theorem~\ref{thm:main}]
Here $n>1$, so the divisors $1,n$ are distinct. Their contributions
to \eqref{eq:def} are $uv^n$ and $vu^n$. For every other divisor,
Lemma~\ref{lem:factorial} makes $Q(n,d)$ divisible by $m!$, hence by $m$.
This proves \eqref{eq:tail}. Apply that identity at $n$ and
$n+\varphi(m)$, then apply Lemma~\ref{lem:powers} to $u$ and $v$.
\end{proof}

\begin{corollary}\label{cor:oeis}
The conjectures in A258899 and A258903 hold. Also,
$c(n)\equiv2\pmod m$ for every $m\geq1$ and $n\geq2m+4$,
proving the conjecture in A121860 (which was posed for $m\geq3$).
\end{corollary}
\begin{proof}
Substitute $(u,v)=(2,1),(2,3),(1,1)$ into Theorem~\ref{thm:main}.
In particular, $a(n)\equiv2+2^n$ and
$b(n)\equiv2\cdot3^n+3\cdot2^n\pmod m$ in the stated range.
An eventual period $T$ implies that the least eventual period divides $T$:
extend the periodic tail to a sequence on $\mathbb Z$; its periods form
a subgroup of $\mathbb Z$, whose positive generator divides $T$.
\end{proof}

\paragraph{A smaller universal period.}
Let $\lambda(m)$ be the exponent of $(\mathbb Z/m\mathbb Z)^\times$,
with $\lambda(1)=1$. Each unit order divides $\varphi(m)$, by partitioning
the unit group into cosets of the cyclic subgroup it generates. Thus their
least common multiple $\lambda(m)$ divides $\varphi(m)$.
In fact \eqref{eq:period} holds with $\lambda(m)$ in place of $\varphi(m)$.
For a prime power $p^e\Vert m$ and $p\nmid z$, the Chinese remainder
theorem lifts $z$ modulo $p^e$ to a unit modulo $m$ (choose residue $1$
at all other prime powers). Therefore $z^{\lambda(m)}\equiv1\pmod{p^e}$.
When $p\mid z$, both shifted powers vanish for $n\geq m$.
This proves the power lemma with $\lambda(m)$, and the same tail identity
finishes the argument. The bound $2m+4$ is a sufficient onset, not a least one.

\section{Companion formulas}
\begin{corollary}\label{cor:egf}
We have $A_{u,v}(n)=A_{v,u}(n)$ and the formal identity
\begin{equation}\label{eq:egf}
 \sum_{n\geq0}\A(n)\frac{x^n}{n!}
 =1+\sum_{d\geq1}\frac{u^d}{d!}\bigl(\exp(vx^d)-1\bigr).
\end{equation}
\end{corollary}
\begin{proof}
The substitution $d\mapsto n/d$ proves symmetry. For $n\geq1$,
the coefficient of $x^n$ on the right is
\[
 \sum_{de=n}\frac{u^dv^e}{d!e!}.
\]
Each coefficient involves finitely many terms.
For the positive weights in the two Hanna entries,
\eqref{eq:egf} equals
\[
 2-\exp(u)+\sum_{d\geq1}u^d\exp(vx^d)/d!
\]
near zero.
This follows also analytically: for $|x|\leq r<1$, the summands are
bounded by $|u|^d\exp(|v|r)/d!$.
Interchanging $u,v$ proves both alternative generating functions in their
FORMULA fields. For A121860, delete the initial $1$ in
\eqref{eq:egf}, matching its offset-one generating function.

\end{proof}
\section{Prime moduli and arbitrary integer weights}
\begin{lemma}\label{lem:primorial}
If $d,e\geq2$ and $n=de$, every prime $p\leq n$ divides $Q(n,d)$.
\end{lemma}
\begin{proof}
Counting multiples of $p,p^2,\ldots$ in each factorial gives
\[
 v_p(Q(n,d))=\sum_{j\geq1}\left(\left\lfloor\frac{de}{p^j}\right\rfloor
 -\left\lfloor\frac d{p^j}\right\rfloor
 -\left\lfloor\frac e{p^j}\right\rfloor\right).
\]
Every summand is nonnegative since $de\geq d+e$.
Interchange $d,e$ so $d\leq e$. Suppose $p\geq3$.
If $e<p$, the $j=1$ summand is positive because $de\geq p$.
If $e\geq p$, it is at least
$(d-1)\lfloor e/p\rfloor-\lfloor d/p\rfloor
\geq d-1-\lfloor d/p\rfloor\geq1$.
For $p=2$ this inequality works when $d\geq3$.
When $d=2,e\geq3$, the $j=1$ summand is
$e-1-\lfloor e/2\rfloor\geq1$; the remaining case has $Q(4,2)=6$.
\end{proof}

\begin{corollary}\label{cor:prime}
For every prime $p$, integer $u,v$, and $n\geq p$,
\[
 \A(n)\equiv uv^n+vu^n\pmod p,\qquad
 \A(n+p-1)\equiv\A(n)\pmod p.
\]
Thus A258899 has period $10$ from $n=11$ modulo $11$,
and A258903 has period $6$ from $n=7$ modulo $7$.
\end{corollary}
\begin{proof}
Lemma~\ref{lem:primorial} removes every proper-divisor term.
For a base not divisible by $p$, the unit-product argument in
Lemma~\ref{lem:powers} gives period $p-1$; a base divisible by $p$
has zero powers for all positive exponents. These are sufficient onsets.
\end{proof}

\begin{corollary}\label{cor:triangle}
For any integer-valued function $W(d,e)$,
$\sum_{d\mid n}W(d,n/d)Q(n,d)$ is congruent to
$W(1,n)+W(n,1)$ modulo $m$ if $n\geq2m+4$, or modulo a prime
$p$ if $n\geq p$. In A386877 \cite{a386877},
all entries strictly between the endpoint $1$'s vanish modulo
the respective modulus in these ranges; the row sum is $c(n)$ for $n\geq1$.
The sign of $T(n,k)$ is the divisibility indicator A113704 for $1\leq k\leq n$.
\end{corollary}
\begin{proof}
Apply Lemma~\ref{lem:factorial} or Lemma~\ref{lem:primorial} termwise.
Nondivisor entries and column zero are zero for $n>0$;
the other entries are the positive integers $Q(n,d)$ by the triangle's definition,
which also proves the sign assertion.
\end{proof}

\section{Direct cross-references and a history observation}
\begin{corollary}\label{cor:interpretations}
The number of matrices with entries $1,\ldots,n$, each used once,
up to independent row and column permutations is $c(n)$.
The same count enumerates the labeled rectangular-band operations
$(i,j)(k,l)=(i,l)$. Also $c(n)=2$ exactly when $n$ is prime.
\end{corollary}
\begin{proof}
For each shape $d\times e$, $de=n$, the action of $S_d\times S_e$
on the $n!$ distinct-entry arrays is free, giving $Q(n,d)$ orbits.
A grid gives the stated operation. Conversely its row and column classes
are recovered from equality of left and right multiplication maps,
so two grids give the same operation exactly when they differ by row
and column permutations. Sum over the shapes.
For prime $n$ only the two endpoint terms occur; for composite $n>1$
a positive proper-divisor term adds to $2$, and $c(1)=1$.
The rectangular-band enumeration and primality criterion already appear in
Devillet and Kiss \cite[Proposition 4.13 and Corollary 4.15]{dk},
whose preprint dates from June 2018.
\end{proof}

\begin{corollary}\label{cor:quet}
For A057625 \cite{a057625}, let $L(n)=\sum_{d\mid n}n!/d!$, $n\geq1$.
Then $L(n)\equiv1\pmod m$ for every $m\geq1$, $n\geq2m$, and
$\sum_{n\geq1}L(n)x^n/n!=\sum_{d\geq1}x^d/(d!(1-x^d))$.
\end{corollary}
\begin{proof}
The term $d=n$ is $1$. For a proper divisor $d\leq n/2$,
$n-d\geq m$, and
$n!/d!=m!(n-d-m)!\binom{n}{d,m,n-d-m}$ is divisible by $m$.
Expand $x^d/(1-x^d)=\sum_{e\geq1}x^{de}$ for the generating function.
\end{proof}

\begin{corollary}\label{cor:other}
Let $P(n),B(n),F(n)$ denote A038041, A061095 and A236696,
respectively \cite{other}; their offset is one. Then
\begin{align*}
 \sum_{n\geq1}P(n)\frac{x^n}{n!}
   &=\sum_{k\geq1}\left(\exp(x^k/k!)-1\right),\\
 \sum_{n\geq1}B(n)\frac{x^n}{n!}
   &=\sum_{k\geq1}\frac{x^k}{k!-x^k},\\
 \sum_{n\geq1}F(n)\frac{x^n}{n!}
   &=\sum_{k\geq1}\left(\exp(k^{k-1}x^k/k!)-1\right).
\end{align*}
Also $P(n)=2$ exactly at prime $n$, and
$0\leq B(n)/n!-1\leq n2^{-n/2}$.
If $H(n)$ is A320444 \cite{a320444}, its prime specialization is
$H(p+1)=1+(p+1)^{p-1}$ for every prime $p$.
\end{corollary}
\begin{proof}
The defining sums are
\begin{align*}
 P(n)&=\sum_{ke=n}\frac{n!}{e!(k!)^e},&
 B(n)&=\sum_{ke=n}\frac{n!}{(k!)^e},\\
 F(n)&=\sum_{ke=n}\frac{n!}{e!}\left(\frac{k^{k-1}}{k!}\right)^e.
\end{align*}
Expand the exponentials and geometric series; every positive-degree
coefficient is a finite sum. The displayed A236696 generating function
omits the subtractions of $1$; its constant coefficient is an infinite
sum of $1$'s. The formula above is the well-defined offset-one version.
For $P$, the endpoint partitions contribute $2$ when $n>1$;
a proper factor supplies additional partitions at composite $n$, and $P(1)=1$.
For $B$, the $k=1$ term is $n!$. If $k\geq2$, then
$k!\geq2^{k-1}\geq2^{k/2}$, so $(k!)^{n/k}\geq2^{n/2}$;
there are at most $n$ other terms, proving the bound and $B(n)\sim n!$.
Finally
$H(n+1)=\sum_{d\mid n}Q(n,d)((n+1)/d)^{n/d-1}$
has only the divisors $1,p$ when $n=p$, contributing
$(p+1)^{p-1}$ and $1$. Thus the printed prime formula in A320444
needs the shift to $H(p+1)$: $H(3)=4$, whereas $17=1+4^2=H(4)$.
These sums have different denominators; no congruence for $P,B,F,H$
is inferred from Theorem~\ref{thm:main}.
\end{proof}

\begin{corollary}\label{cor:three}
Three consecutive residues $2$ in A121860 need not begin a constant tail.
Modulo $89$, $(c(41),c(42),c(43),c(44))\equiv(2,2,2,76)$.
\end{corollary}
\begin{proof}
The numbers $41,43$ are prime. Pairing complementary proper divisors gives
$c(42)=2+2(Q(42,2)+Q(42,3)+Q(42,6))$ and
$c(44)=2+2(Q(44,2)+Q(44,4))$.
The quotient residues modulo $89$ are $(82,71,25)$ and $(21,16)$,
respectively, giving the asserted residues by exact arithmetic.
This corrects David Pasino's observation in the A121860 history,
15 November 2025, discussion following revision 36 \cite{history}.
\end{proof}

\begingroup\small
\begin{longtable}{@{}>{\raggedright\arraybackslash}p{0.25\textwidth}>{\raggedright\arraybackslash}p{0.39\textwidth}>{\raggedright\arraybackslash}p{0.29\textwidth}@{}}
\toprule
Source item & Statement & Coverage \\
\midrule
\endhead
A121860, COMMENTS (2019, 2025) & Matrix and rectangular-band counts & Proved; prior band proof cited: Corollary~\ref{cor:interpretations} \\
A121860 history, 15 Nov 2025 & Three residues $2$ predict onset & Refuted-and-corrected: Corollary~\ref{cor:three} \\
A386877 & Row sum and vanishing interior entries & Proved: Corollary~\ref{cor:triangle} \\
A057625 & Constant tail $1$; generating function & Proved: Corollary~\ref{cor:quet} \\
A038041 & Prime criterion; generating function & Proved: Corollary~\ref{cor:other} \\
A061095 & Generating function; $B(n)\sim n!$ & Proved: Corollary~\ref{cor:other} \\
A236696, FORMULA & Generating function's constant coefficient & Refuted-and-corrected: Corollary~\ref{cor:other} \\
A320444, FORMULA & Prime-specialization index & Refuted-and-corrected: Corollary~\ref{cor:other} \\
A100195 \cite{a100195} & Record Bernoulli denominators & Not covered: different arithmetic question \\
A038041, A061095, A236696, A320444 & Possible modular tails & Not covered: denominators differ \\
\bottomrule
\end{longtable}
\endgroup


\section{Coverage of the posing assertions}
\begin{center}\small
\begin{tabular}{@{}>{\raggedright\arraybackslash}p{0.23\textwidth}>{\raggedright\arraybackslash}p{0.39\textwidth}>{\raggedright\arraybackslash}p{0.29\textwidth}@{}}
\toprule
Source item & Statement & Coverage \\
\midrule
A258899, COMMENTS, 3 Aug 2025 & Totient-dividing eventual period & Theorem~\ref{thm:main}, $(u,v)=(2,1)$ \\
A258903, COMMENTS, 8 Aug 2025 & Totient-dividing eventual period & Theorem~\ref{thm:main}, $(u,v)=(2,3)$ \\
A121860, COMMENTS, 8 Aug 2025 & Constant residue $2$ for $m\geq3$ & Corollary~\ref{cor:oeis}; also $m=1,2$ \\
Both Hanna entries, FORMULA & Alternative exponential generating functions & Equation~\eqref{eq:egf} and symmetry \\
A121860, first COMMENT & $c(n)=2$ exactly for prime $n$ & Prior proof cited: \cite[Corollary 4.15]{dk}; Corollary~\ref{cor:interpretations} \\
Entry examples of early onset & Onsets $11$ and $7$ at moduli $11$ and $7$ & Proved for all later indices: Corollary~\ref{cor:prime} \\
\bottomrule
\end{tabular}
\end{center}

\begingroup\raggedright
\paragraph{Lean declarations.}
In the pinned
\href{https://github.com/the-omega-institute/trureturing/blob/eb2703740e8b4b97e408a80b2eb0d0940939dfc3/D5/S1/Recurrence/Periodic/BalaDivisorFactorialTotientPeriod.lean}{\texttt{BalaDivisorFactorialTotientPeriod.lean}},
\path{D5.S1.Recurrence.Periodic.BalaDivisorFactorialTotientPeriod.result}
proves the two natural-weight sequence statements, with
$N=2m+4$ and $T=\varphi(m)$.
The helper \texttt{proper\_factorial\_quotient\_dvd} proves divisibility
by $m$, and \texttt{divisor\_sum\_tail} proves \eqref{eq:tail} for
natural $u,v$. The imported
\href{https://github.com/the-omega-institute/trureturing/blob/eb2703740e8b4b97e408a80b2eb0d0940939dfc3/D5/S1/Recurrence/Periodic/StirlingTransformTotientPeriod.lean}{power-period Lean file}
contains \texttt{pow\_add\_totient}, the natural-base power lemma with
$n\geq m$. The signed-weight extension, the $m!$ strengthening,
the A121860 application, the Carmichael period, and the generating-function
identities have the written proofs above, as do the prime-modulus bound,
arbitrary integer weights, and the cross-referenced corollaries.
The accompanying \texttt{verify.py} supplies exact finite checks of the
divisor sums, the power and tail congruences, and the listed sequence values;
its output states the finite ranges explicitly.\par
\endgroup

\begin{thebibliography}{10}
\bibitem{a258899}
P. D. Hanna, OEIS A258899, entry dated 20 June 2015;
P. Bala, COMMENTS, 3 August 2025, and FORMULA, 4 August 2025.
\url{https://oeis.org/A258899}. Accessed 3 November 2026.
\bibitem{a258903}
P. D. Hanna, OEIS A258903, entry dated 20 June 2015;
P. Bala, COMMENTS and FORMULA, 8 August 2025.
\url{https://oeis.org/A258903}. Accessed 3 November 2026.
\bibitem{a121860}
V. Jovovic, OEIS A121860, entry dated 9 September 2006;
P. Bala, COMMENTS, 8 August 2025.
\url{https://oeis.org/A121860}. Accessed 3 November 2026.
\bibitem{dk}
J. Devillet and G. Kiss, \emph{Characterizations of biselective operations},
arXiv:1806.02073v1, 6 June 2018; Proposition 4.13 and Corollary 4.15.
\url{https://arxiv.org/abs/1806.02073v1}.
\bibitem{a386877}
P. Bala and P. Luschny, OEIS A386877, 9 August 2025.
\url{https://oeis.org/A386877}.
\bibitem{a057625}
L. Quet, OEIS A057625, 9 October 2000.
\url{https://oeis.org/A057625}.
\bibitem{other}
C. G. Bower, OEIS A038041; H. Bottomley, A061095, 29 May 2001;
E. Munarini, A236696, 30 January 2014.
\url{https://oeis.org/A038041}, \url{https://oeis.org/A061095},
\url{https://oeis.org/A236696}. FORMULA fields.
\bibitem{a320444}
G. Wiseman, OEIS A320444, 9 January 2019. FORMULA field.
\url{https://oeis.org/A320444}.
\bibitem{history}
D. Pasino, discussion following revision 36 of OEIS A121860,
15 November 2025. \url{https://oeis.org/history?seq=A121860}.
\bibitem{a100195}
E. W. Weisstein, OEIS A100195, 8 November 2004.
\url{https://oeis.org/A100195}.
\end{thebibliography}
\medskip\noindent\small
\copyright\ 2026 Haobo Ma and Wenlin Zhang.
Licensed under Creative Commons Attribution 4.0 International (CC BY 4.0),
\url{https://creativecommons.org/licenses/by/4.0/}.
\end{document}
